\begin{abstract}
Training multilingual language models requires repeated decisions about data mixtures and architectures, often based on evaluations of small proxy models.
Yet these evaluations are useful only if results from small-scale models follow the same trends as their larger-scale counterparts, a challenge compounded in multilingual settings where many language–benchmark pairs exhibit little measurable signal at small scales.
This work asks four foundational questions around this practice: (1) for multilingual language models, at what model size do evaluations become sufficiently predictive of larger models? (2) how early during pretraining can decisions be made? (3) can we determine before doing many expensive evaluations whether a particular evaluation is likely to be informative at small-scale? and (4) how do these patterns change as multilingual training data varies in composition? 
We pretrain a controlled ladder of 162 models spanning five sizes from 175M to 1.7B non-embedding parameters and seven language settings from 1 to 100 languages, with four design interventions and seed replicates. 
We evaluate every checkpoint on per-language bits per byte in 100 languages and task performance on 14 multilingual benchmark families.
% RQ1
We characterize which benchmark families exhibit predictable scaling across model size and training, and distinguish benchmarks whose performance improves consistently from those whose scaling is weak, noisy, or non-monotonic. 
% RQ2
For each benchmark, we quantify how well performance at smaller model sizes predicts performance at larger sizes, and how early in training this correspondence emerges. We then test whether these small-scale evaluations recover the relative effects of our data-mixture and architecture interventions observed at full scale.
% RQ3
Finally, we ask whether the reliability of a small-scale evaluation can itself be anticipated cheaply, comparing measurements that assess whether a benchmark has sufficient signal to recover full-scale results across model scales and levels of multilingual coverage. 
% RQ4
% Finally, we test whether scaling behavior learned from observed languages transfers to held-out languages, and how this transfer depends on the underlying design intervention. 
% We find that likelihood-based measurements scale predictably while four-option knowledge benchmarks sit at chance up to 1.7B; that a 350M proxy at 20\% of its run already reads a data decision the way the 1.7B reference does on bits per byte, while the benchmark suite reads every decision as a coin flip; that the reliability of a proxy is set by the size of the intervention's effect relative to seed noise rather than by the intervention itself; and that scaling exponents transfer across languages, so that one small measurement predicts a never-trained language's bits per byte at 1.7B within 7\%. 
We release the ladder and the per-checkpoint measurements to support cheaper, better-informed multilingual model development.
\end{abstract}
